\documentclass{article}
\usepackage{amsmath,amssymb,amsfonts}
\usepackage[a4paper,margin=1in]{geometry}

\title{Alpha–Beta–Gamma 图表公式}
\author{}
\date{}

\usepackage[UTF8]{ctex}
\begin{document}
\maketitle

\section*{符号与数据}
设$G=(V,E)$为一个无向图，其中$V$为顶点集，$E\subseteq\{\{i,j\}:i,j\in V,\ i\neq j\}$为边集。设$C$为颜色集合，而$V^c\subseteq V$表示具有颜色$c\in C$的顶点集合。对于每一个$i\in V$，令
\[
N(i)=\{j\in V:\{i,j\}\in E\}
\]
表示$i$的邻域。对于每一个$U\subseteq V$，定义由此产生的边集和割集。
\[
E(U)=\{\{i,j\}\in E:i\in U,\ j\in U\},
\qquad
\delta(U)=\{\{i,j\}\in E:\lvert\{i,j\}\cap U\rvert=1\}.
\]
对于每一个$W\subseteq V$，让$t_W$表示所诱导的图$G[W]$的给定边连通性参数。涉及$1/|N(i)|$和$1/|V^c|$的约束条件，是针对带有$|N(i)|>0$和$|V^c|>0$的索引来定义的。

\section*{决策变量}
对于每个$i\in V$，定义二进制变量。
\[
\alpha_i=\begin{cases}
1,&\text{if vertex $i$ belongs to the selected component }V_B,\\
0,&\text{otherwise,}
\end{cases}
\qquad
\beta_i=\begin{cases}
1,&\text{if vertex $i$ belongs to the boundary-neighbor set }H,\\
0,&\text{otherwise.}
\end{cases}
\]
对于每个$\{i,j\}\in E$，定义二元变量$\gamma_{ij}$，其中$\gamma_{ij}=1$表示边$\{i,j\}$属于两个选定的生成树中的一个。

\section*{数学公式如下：}
\begin{align}
\max_{\alpha,\beta,\gamma}\quad
&\sum_{\{i,j\}\in E}\left(\alpha_i\beta_j+\alpha_j\beta_i\right) \\[-1mm]
\text{s.t.}\quad
&\sum_{i\in V}\alpha_i\ge 1, \label{eq:nonempty-alpha}\\
&\sum_{i\in V}\beta_i\ge 1, \label{eq:nonempty-beta}\\
&\beta_i\le 1-\alpha_i &&\forall i\in V, \label{eq:beta-disjoint}\\
&\beta_i\le \sum_{j\in N(i)}\alpha_j &&\forall i\in V, \label{eq:beta-adjacent}\\
&\beta_i\ge \frac{1}{|N(i)|}\sum_{j\in N(i)}\alpha_j-\alpha_i &&\forall i\in V, \label{eq:beta-forced}\\
&\frac{1}{|V^c|}\sum_{i\in V^c}\beta_i
  \le \sum_{i\in V^c}\alpha_i\le 1 &&\forall c\in C, \label{eq:color}\\
&\sum_{\{i,j\}\in E}\gamma_{ij}=|V|-2, \label{eq:tree-count}\\
&\sum_{\{i,j\}\in E(U)}\gamma_{ij}\le |U|-1 &&\forall \varnothing\ne U\subseteq V, \label{eq:acyclicity}\\
&\alpha_i+\alpha_j-2\alpha_i\alpha_j\le 1-\gamma_{ij}
 &&\forall\{i,j\}\in E, \label{eq:same-component}\\
&\sum_{i\notin W}\alpha_i+\sum_{i\in W}(1-\alpha_i)\ge 1
 &&\forall W\subseteq V:\ t_W\le |\delta(W)|, \label{eq:nogood}\\
&\alpha_i\in\{0,1\} &&\forall i\in V, \label{eq:alpha-domain}\\
&\beta_i\in\{0,1\} &&\forall i\in V, \label{eq:beta-domain}\\
&\gamma_{ij}\in\{0,1\} &&\forall\{i,j\}\in E. \label{eq:gamma-domain}
\end{align}

\end{document}
